\section{What Pontifical Right Does---and Does Not---Mean}

The 18 October decree placed the FSSP within the Code's law of societies of apostolic life. Canon 731 describes bodies resembling institutes of consecrated life whose members, without religious vows as such, pursue the society's apostolic purpose and strive for charity through a common life governed by constitutions. The FSSP is clerical: priestly formation and ministry define its purpose, and its clerics are ordinarily incardinated in the society under canon 736 and its proper law.

The term ``pontifical right'' describes the level of ecclesiastical erection and supervision. The erection decree subjected the Fraternity to the Roman Pontiff through the Pontifical Commission \emph{Ecclesia Dei} until otherwise provided. It did not turn each house into an enclave outside a diocese. Canon 733 requires the previous written consent of the diocesan bishop before the society erects a house and establishes a local community. Canon 738 subjects members to their own moderators in internal life and discipline, and also to the diocesan bishop in public worship, care of souls, and apostolic works.

\begin{threestable}{Relation}{Competent authority}{Historical consequence}
Internal life, formation, incorporation, assignments, and discipline & Superior general and subordinate superiors under constitutions and Holy See supervision & The FSSP can preserve common formation and patrimony across dioceses.\\
Erection of a house & Society authority after prior written diocesan consent & A world map of FSSP houses is also a history of episcopal receptions, not autonomous occupation.\\
Public worship, care of souls, and diocesan apostolate & Diocesan bishop, coordinated with FSSP superiors and any written agreement & Pontifical right does not make local consent or diocesan governance optional.\\
Proper liturgical faculty & Holy See acts, proper law, and the conditions stated in each act & The faculty can be corporate and stable while still carrying location and consent conditions.\\
\end{threestable}

This divided competence is not a defect accidentally imposed on the FSSP. It is how a transdiocesan society performs local ministry. General government protects continuity of formation and personnel; a bishop remains responsible for worship and souls in his particular church. The 1988 decree underscored that members should cultivate communion with the bishop and diocesan presbyterate and observe the law governing penance, marriage, and parish registers.

\subsection{Incardination, incorporation, and association}

The FSSP's current statistics separate incardinated priests, priests incorporated \emph{ad annum}, postulants, and associated priests. Those are institutional categories, not synonyms. Incardination is the canonical bond by which a cleric belongs to a particular church or qualified institute or society; incorporation describes membership under the society's proper law. A priest can be associated with the work without becoming incardinated. A seminarian or deacon is a member in formation but is not already counted among priests.

This distinction matters historically because the FSSP was created to reproduce a clergy, not merely to coordinate celebrants. It needed a recognized process of admission and incorporation, seminaries, bishops willing and authorized to confer orders, superiors able to assign priests, and communities able to sustain common life. Growth in Mass centers is therefore only one measure of institutional maturity.

\subsection{Parish, chaplaincy, and Mass center}

A personal parish erected under canon 518 supplies a stable pastoral structure defined by rite, language, nationality, or another reason rather than territory alone. A chaplaincy, oratory, entrusted church, and recurring Mass center can each provide serious pastoral life without being a personal parish. The FSSP's 2025 report counts 48 personal parishes and 251 Mass centers, confirming that most centers do not possess parish status.

The Ottawa example shows why the date and category must be named. The older-liturgical community predated the FSSP. FSSP administration began in 1995, and the archbishop erected the personal parish in 1997. Calling it simply a parish ``founded by the FSSP'' would erase nearly three decades of local history; calling every FSSP site a parish would erase the bishop's juridical act.

\claimcheck{``Pontifical right means the local bishop cannot restrict an FSSP apostolate.''}{Pontifical right protects a real relation to the Holy See and rightful autonomy in internal life. Canons 733 and 738 preserve the diocesan bishop's consent and governance in the local apostolate. The scope of a particular liturgical decree must also be read from its own conditions.}
